\section{Controlled Downstream Check}\label{app:downstream_temporal_invalidity}

This appendix records the controlled synthetic downstream benchmark used to
probe the predictive cost of temporal invalidity. The setup is intentionally
minimal: a phasewise drifting binary classification stream, a windowed logistic
backend, and four policies comparing long memory, short memory, ADWIN-driven
memory, and UMR-regulated memory. The comparison is restricted to the phase-2
cap-only interval, where UMR has already contracted but ADWIN remains silent.

\begin{figure}[!htbp]
\centering
\safeincludegraphics[width=0.95\linewidth]{figures/temporal_invalidity/fig_temporal_invalidity}
\caption{Controlled downstream benchmark. Top: true drift rate. Middle:
effective horizons and the cap-only interval. Bottom: rolling prequential
accuracy. The representative seed shows a 3.12 point gain for UMR over
Fixed-400 within the cap-only window.}
\label{fig:temporal_invalidity}
\end{figure}

\begin{table}[!htbp]
\centering
\small
\caption{Controlled downstream benchmark summary. The reported delta is UMR
minus Fixed-400 within the cap-only interval, averaged across six seeds.}
\label{tab:temporal_invalidity}
\begin{tabular}{lrrrr}
\toprule
Method & Global acc. & Cap-only acc. & $\Delta$ cap-only & Mean cap-only width \\
\midrule
Fixed-400 & 0.9304 & 0.9040 & 0.00 & 293.3 \\
Fixed-100 & 0.9598 & 0.9576 & +5.36 & 293.3 \\
ADWIN & 0.9340 & 0.9040 & 0.00 & 293.3 \\
UMR & 0.9415 & 0.9205 & +1.66 & 293.3 \\
\bottomrule
\end{tabular}
\end{table}

The result is consistent in direction across all six seeds, but the effect is
modest in magnitude. We therefore treat it as mechanistic supporting evidence
rather than as a central claim.
